\documentclass[12pt]{article}

\usepackage{amsmath,amssymb,amsthm,mathtools,enumitem}

\frenchspacing
\allowdisplaybreaks[2]

\newtheorem{theorem}{Theorem}[section]
\newtheorem{lemma}[theorem]{Lemma}
\newtheorem{proposition}[theorem]{Proposition}
\newtheorem{corollary}[theorem]{Corollary}
\theoremstyle{remark}
\newtheorem{remark}[theorem]{Remark}

\DeclareMathOperator{\E}{\mathbb E}
\DeclareMathOperator{\PP}{\mathbb P}
\DeclareMathOperator{\R}{\mathbb R}
\DeclareMathOperator{\Z}{\mathbb Z}
\DeclareMathOperator{\N}{\mathbb N}
\DeclareMathOperator{\T}{\mathbb T}
\DeclareMathOperator{\divg}{div}
\DeclareMathOperator{\supp}{supp}
\DeclareMathOperator{\KL}{KL}
\DeclareMathOperator{\TV}{TV}

\newcommand{\D}{D}
\newcommand{\one}{\mathbf 1}
\newcommand{\norm}[1]{\left\lVert #1\right\rVert}
\newcommand{\abs}[1]{\left\lvert #1\right\rvert}
\newcommand{\asym}{\asymp}

\title{Anisotropic Linearization and a Frozen Minimax Exponent\\
for Darcy Coefficient Recovery}
\author{Open-weight synthesis article}
\date{}

\begin{document}

\maketitle

\begin{abstract}
We study the statistical inverse problem of recovering a Darcy conductivity \(f\) from Gaussian white-noise observations of the potential \(u_f\) solving \(\divg(f\nabla u)=1\) with zero boundary data on the unit ball \(D\subset \R^d\), \(d\ge 2\). We prove the linearization at the constant conductivity, including the exact radial identity and the anisotropic frozen symbol. We show that the natural fixed-anchor quadratic remainder estimate at \(f\equiv 1\) is false and prove the correct moving-anchor quadratic estimate. For a frozen periodic linearized model with nonzero frozen direction, we prove the minimax \(L^2\) exponent
\[
n^{-2\alpha/(2\alpha+d+3)}
\]
(up to logarithmic factors in the certified statement; the periodic wavelet construction below gives the sharper constant-order version). We also prove that the phase-space variance of the genuine variable-coefficient linearized symbol has order \(R^{d+3}\), and we prove a rigorous two-point lower bound at the critical point \(x=0\) giving the weaker exponent
\[
n^{-(2\alpha+d)/(2\alpha+d+4)}.
\]
The full nonlinear variable-coefficient minimax rate over the original Darcy class remains open.
\end{abstract}

\section{Introduction and certified scope}
\label{sec:intro}

Let
\[
D=\{x\in \R^d: |x|<1\},\qquad d\ge 2.
\]
For a uniformly positive conductivity \(f\), let \(u_f=\Phi(f)\) denote the solution of Darcy's problem
\begin{equation}
\label{eq:darcy}
\divg(f\nabla u)=1 \quad \text{in }D,\qquad u|_{\partial D}=0.
\end{equation}
The observation model is Gaussian white noise:
\[
Y_n = u_f + \frac1{\sqrt n}\,\mathbb W,
\]
where \(\mathbb W\) is Gaussian white noise on \(L^2(D)\). The parameter class is
\[
\mathcal F(\alpha,r)
=
\left\{
h\in C^\alpha(\overline D): \norm{h-\one}_{C^\alpha(\overline D)}\le r
\right\},
\]
with \(\alpha>1\) and fixed \(r\in(0,1)\). The original minimax question is the asymptotic size, up to logarithmic factors, of
\[
\inf_{\widehat f_n}\sup_{f_0\in\mathcal F(\alpha,r)}
\E_{f_0}\left[
\norm{\widehat f_n-f_0}_{L^2(D)}^2
\right].
\]

This paper proves a scoped set of results that identify the linearized anisotropy and the resulting frozen minimax exponent, while leaving the full nonlinear variable-coefficient problem open. The main certified results are:
\begin{enumerate}[label=(\roman*)]
\item the linearization at \(f\equiv1\), the exact radial identity, and injectivity of the linearized operator on regular functions;
\item the failure of the fixed-anchor quadratic estimate at \(f\equiv1\), and the validity of the moving-anchor quadratic estimate;
\item the minimax rate of a frozen periodic linearized model with nonzero frozen direction, giving the exponent \(2\alpha/(2\alpha+d+3)\);
\item a phase-space variance calculation for the variable-coefficient linearized symbol, also of order \(R^{d+3}\);
\item a critical-point two-point lower bound for the original nonlinear problem, giving the weaker exponent \((2\alpha+d)/(2\alpha+d+4)\).
\end{enumerate}
The frozen exponent is the candidate exponent suggested by the anisotropic linearization. The critical-point lower bound is strictly weaker and shows that the origin, although more degenerate locally, is not the minimax bottleneck.

Throughout, \(A\asym B\) means \(cA\le B\le CA\), with constants depending only on the structural parameters indicated in the statement. Constants denoted by \(C,c\) may change from line to line.

\section{Preliminaries}
\label{sec:prelim}

We use the weak formulation of \eqref{eq:darcy}. For \(f\in L^\infty(D)\) with \(f\ge \lambda>0\), \(u_f\in H^1_0(D)\) is the unique solution of
\begin{equation}
\label{eq:weak-forward}
\int_D f\nabla u_f\cdot \nabla\varphi\,dx
=
-\int_D \varphi\,dx
\qquad \forall \varphi\in H^1_0(D).
\end{equation}
Let \(C_P\) denote the Poincar\'e constant of \(D\).

% RECORDS: 129,756
\begin{lemma}[Forward well-posedness and uniform energy bound]
\label{lem:forward-energy}
Let \(f\in L^\infty(D)\) and \(f\ge \lambda>0\). Then \eqref{eq:weak-forward} has a unique solution \(u_f\in H^1_0(D)\), and
\[
\norm{\nabla u_f}_{L^2(D)}\le \frac{C_P}{\lambda},
\qquad
\norm{u_f}_{L^2(D)}\le \frac{C_P^2}{\lambda}.
\]
\end{lemma}
\begin{proof}
The bilinear form
\[
B(u,\varphi)=\int_D f\nabla u\cdot\nabla\varphi\,dx
\]
is continuous and coercive on \(H^1_0(D)\), since
\[
B(\varphi,\varphi)\ge \lambda\norm{\nabla\varphi}_{L^2}^2.
\]
The right-hand side \(\varphi\mapsto -\int_D\varphi\) is bounded on \(H^1_0(D)\) by Poincar\'e's inequality. Existence and uniqueness follow from the Lax--Milgram theorem; see \cite{Evans}.

Testing \eqref{eq:weak-forward} with \(\varphi=u_f\) gives
\[
\int_D f|\nabla u_f|^2\,dx
=
-\int_D u_f\,dx.
\]
Hence
\[
\lambda\norm{\nabla u_f}_2^2
\le
\norm{u_f}_2
\le
C_P\norm{\nabla u_f}_2,
\]
so \(\norm{\nabla u_f}_2\le C_P/\lambda\). A second application of Poincar\'e gives the \(L^2\) bound.
\end{proof}

For \(f\ge\lambda>0\) and \(h\in L^\infty(D)\), define the linearized operator \(A_f h\) weakly by
\begin{equation}
\label{eq:Af-def}
\int_D f\nabla(A_f h)\cdot\nabla\varphi\,dx
=
-\int_D h\nabla u_f\cdot\nabla\varphi\,dx
\qquad \forall\varphi\in H^1_0(D).
\end{equation}
Again by Lax--Milgram, \(A_f h\in H^1_0(D)\) is well-defined whenever \(h\nabla u_f\in L^2(D)\); in particular this holds for \(h\in L^\infty(D)\). On the regular positive cone this operator is the Fr\'echet derivative of the forward map \(\Phi\).

We shall use the standard Cameron--Martin formula for Gaussian white noise. If \(P_s\) denotes the law of \(s+n^{-1/2}\mathbb W\), then for \(s,t\in L^2\),
\[
\KL(P_s\|P_t)=\frac n2\norm{s-t}_{L^2}^2.
\]
For completeness we record this as a lemma.

% RECORDS: 199,756
\begin{lemma}[KL divergence in Gaussian white noise]
\label{lem:white-kl}
Let \(\{e_k\}_{k\ge1}\) be an orthonormal basis of \(L^2(D)\). Under the white-noise model, the coordinates
\[
Y_k=\langle Y_n,e_k\rangle
\]
are independent Gaussian random variables with variance \(1/n\) and means \(\langle s,e_k\rangle\). Therefore, for \(s,t\in L^2(D)\),
\[
\KL(P_s\|P_t)
=
\sum_{k\ge1}\frac{n}{2}\bigl(\langle s-t,e_k\rangle\bigr)^2
=
\frac n2\norm{s-t}_{L^2(D)}^2.
\]
\end{lemma}
\begin{proof}
For one coordinate, \(Y_k\sim N(\langle s,e_k\rangle,1/n)\) under \(P_s\) and \(Y_k\sim N(\langle t,e_k\rangle,1/n)\) under \(P_t\). The KL divergence between two Gaussian laws with equal variance \(\sigma^2\) is \((\mu-\nu)^2/(2\sigma^2)\). Here \(\sigma^2=1/n\), hence the \(k\)-th contribution is
\[
\frac n2\bigl(\langle s,e_k\rangle-\langle t,e_k\rangle\bigr)^2.
\]
Summing over \(k\) and applying Parseval gives the claim.
\end{proof}

For H\"older--Zygmund spaces we use the standard norm \(C^\alpha(\overline D)\), including the supremum norm. For non-integer \(\alpha\), \(C^\alpha\) coincides with the Besov space \(B^\alpha_{\infty,\infty}\); see \cite{Triebel}. In the frozen periodic model we use the Besov formulation directly.

\section{Linearization at the constant conductivity}
\label{sec:linearization}

% RECORDS: 44,199
\begin{proposition}[Forward solution and derivative at \(f=1\)]
\label{prop:linearization-one}
Let
\[
u_1(x)=\frac{|x|^2-1}{2d}.
\]
Then \(\Phi(1)=u_1\), \(\Delta u_1=1\), \(\nabla u_1=x/d\), and \(u_1|_{\partial D}=0\). For smooth \(h\), the Fr\'echet derivative \(A=A_1=\Phi'(1)\) is given by
\[
Ah=w,
\qquad
\Delta w=-\left(\frac{x}{d}\cdot\nabla h+h\right),
\qquad
w|_{\partial D}=0.
\]
Equivalently,
\[
-\Delta w=\frac{x}{d}\cdot\nabla h+h.
\]
If \(h\equiv g\) is constant, then
\[
Ag=-g u_1.
\]
\end{proposition}
\begin{proof}
A direct calculation gives
\[
\Delta u_1=\frac{1}{2d}\Delta(|x|^2)=1,
\]
and \(u_1=0\) on \(|x|=1\). Hence \(u_1\) solves \eqref{eq:darcy} for \(f=1\).

Let \(f_\varepsilon=1+\varepsilon h\) and write
\[
u_{f_\varepsilon}=u_1+\varepsilon w+o(\varepsilon)
\]
in the weak sense. The weak equation is
\[
\int_D (1+\varepsilon h)\nabla u_{f_\varepsilon}\cdot\nabla\varphi\,dx
=
-\int_D\varphi\,dx.
\]
Taking the first-order term in \(\varepsilon\) gives
\[
\int_D \nabla w\cdot\nabla\varphi\,dx
+
\int_D h\nabla u_1\cdot\nabla\varphi\,dx
=0.
\]
Since \(\nabla u_1=x/d\), this is the weak form of
\[
-\Delta w=\divg\left(h\frac{x}{d}\right)
=
\frac{x}{d}\cdot\nabla h+h.
\]
Equivalently,
\[
\Delta w=-\left(\frac{x}{d}\cdot\nabla h+h\right).
\]

If \(h\equiv g\), then the right-hand side is \(-g\). Since \(\Delta u_1=1\) and \(u_1|_{\partial D}=0\), the unique solution is \(w=-g u_1\).
\end{proof}

% RECORDS: 44
\begin{proposition}[Exact radial identity]
\label{prop:radial-identity}
Let \(h(x)=\phi(|x|)\) be radial and sufficiently smooth. Then \(w=Ah\) is radial. Writing \(w(x)=W(\rho)\), \(\rho=|x|\), one has
\[
W'(\rho)=-\frac{\rho}{d}\phi(\rho),
\qquad 0<\rho<1.
\]
Equivalently,
\[
W(\rho)=\int_\rho^1 \frac{s}{d}\phi(s)\,ds.
\]
\end{proposition}
\begin{proof}
By rotational invariance and uniqueness, \(w\) is radial. The radial Laplacian is
\[
\Delta w=W''(\rho)+\frac{d-1}{\rho}W'(\rho).
\]
For \(h(x)=\phi(\rho)\),
\[
\frac{x}{d}\cdot\nabla h+h
=
\frac{\rho}{d}\phi'(\rho)+\phi(\rho).
\]
Thus \(W\) satisfies
\[
W''+\frac{d-1}{\rho}W'
=
-\frac{\rho}{d}\phi'-\phi.
\]
If \(W'(\rho)=-(\rho/d)\phi(\rho)\), then
\[
W''=-\frac1d\phi-\frac{\rho}{d}\phi',
\]
and
\[
\frac{d-1}{\rho}W'
=
-\frac{d-1}{d}\phi.
\]
Adding gives
\[
-\frac1d\phi-\frac{\rho}{d}\phi'-\frac{d-1}{d}\phi
=
-\phi-\frac{\rho}{d}\phi',
\]
which is exactly the radial equation. Regularity at \(\rho=0\) removes the singular integration constant, and \(W(1)=0\) yields the integral representation.
\end{proof}

% RECORDS: 44,199
\begin{proposition}[Injectivity of \(A\) on regular functions]
\label{prop:injectivity}
If \(h\in C^\alpha(\overline D)\), \(\alpha>1\), and \(Ah=0\), then \(h=0\).
\end{proposition}
\begin{proof}
Since the Dirichlet Laplacian is injective, \(Ah=0\) implies
\[
\frac{x}{d}\cdot\nabla h+h=0
\]
in the distributional sense. Because \(\alpha>1\), \(h\in C^1\), so the identity holds classically. In polar coordinates \(x=\rho\omega\),
\[
\frac{\rho}{d}\partial_\rho h(\rho\omega)+h(\rho\omega)=0,
\]
or equivalently
\[
\partial_\rho\bigl(\rho^d h(\rho\omega)\bigr)=0.
\]
Thus
\[
h(\rho\omega)=\rho^{-d}g(\omega).
\]
Since \(h\) is bounded at \(\rho=0\), necessarily \(g(\omega)=0\) for every \(\omega\). Hence \(h=0\).
\end{proof}

\begin{remark}
Formally freezing the linearized operator at a nonzero point \(x_0\) with \(v=x_0/d\neq0\) gives the constant-coefficient Fourier multiplier
\[
a_v(\xi)=\frac{1+i\xi\cdot v}{|\xi|^2}.
\]
The radial identity shows degree-one smoothing in radial directions, while the frozen symbol exhibits an equatorial degeneracy where \(\xi\cdot v\) is small. Finite-difference computations corroborate the tangential smoothing but are not used in the proofs.
\end{remark}

\section{Anchor correction and moving-anchor stability}
\label{sec:anchor}

The natural fixed-anchor estimate with \(A=A_1\) would assert that
\[
\norm{\Phi(f)-\Phi(f')-A_1(f-f')}_{L^2(D)}
\le
C\norm{f-f'}_{C^\alpha}^2
\]
uniformly on \(\mathcal F(\alpha,r)\). This is false.

% RECORDS: 129
\begin{proposition}[Failure of the fixed-anchor quadratic estimate]
\label{prop:fixed-anchor-false}
Fix \(r\in(0,1)\). There is no finite constant \(C\) such that
\[
\norm{\Phi(f)-\Phi(f')-A_1(f-f')}_{L^2(D)}
\le
C\norm{f-f'}_{C^\alpha}^2
\]
holds for all \(f,f'\in\mathcal F(\alpha,r)\).
\end{proposition}
\begin{proof}
It suffices to use constant conductivities. Let
\[
b=1+\frac r2,
\qquad
0<g<\frac r2,
\]
and set
\[
f=b,
\qquad
f'=b-g.
\]
Both belong to \(\mathcal F(\alpha,r)\), because their deviations from \(1\) are at most \(r/2\). For a constant conductivity \(c>0\), uniqueness and linearity give
\[
\Phi(c)=\frac{u_1}{c}.
\]
By Proposition~\ref{prop:linearization-one}, for the constant increment \(g\),
\[
A_1 g=-g u_1.
\]
Therefore
\[
\Phi(f)-\Phi(f')-A_1(f-f')
=
u_1\left(
\frac1b-\frac1{b-g}+g
\right).
\]
But
\[
\frac1b-\frac1{b-g}+g
=
g\left(1-\frac1{b(b-g)}\right).
\]
As \(g\downarrow0\),
\[
1-\frac1{b(b-g)}
\longrightarrow
1-\frac1{b^2}>0,
\]
because \(b>1\). Since \(\norm{f-f'}_{C^\alpha}=g\) for constant increments,
\[
\frac{
\norm{\Phi(f)-\Phi(f')-A_1(f-f')}_{L^2}
}{
\norm{f-f'}_{C^\alpha}^2
}
\ge
\frac{c}{g}\to\infty.
\]
Thus no finite constant \(C\) can exist.
\end{proof}

The correct estimate anchors the linearization at the base point \(f'\).

% RECORDS: 129
\begin{theorem}[Uniform moving-anchor quadratic remainder]
\label{thm:moving-anchor}
Let \(\alpha>1\) and \(r\in(0,1)\). There exists \(C=C(d,\alpha,r)<\infty\) such that for all \(f,f'\in\mathcal F(\alpha,r)\),
\[
\norm{\Phi(f)-\Phi(f')-A_{f'}(f-f')}_{L^2(D)}
\le
C\norm{f-f'}_{C^\alpha(\overline D)}^2.
\]
One may take \(C\le C_d(1-r)^{-3}\); in particular the constant blows up only as \(r\uparrow1\).
\end{theorem}
\begin{proof}
Write
\[
g=f-f',
\qquad
u=\Phi(f),
\qquad
u'=\Phi(f'),
\qquad
v=A_{f'}g,
\qquad
R=u-u'-v.
\]
Since \(f,f'\in\mathcal F(\alpha,r)\),
\[
f,f'\ge 1-r=:\lambda>0,
\qquad
\norm{g}_{L^\infty}\le \norm{g}_{C^\alpha}.
\]
Let \(q=u-u'\). From the weak equations for \(u\) and \(u'\),
\[
\int_D f'\nabla q\cdot\nabla\varphi\,dx
=
-\int_D g\nabla u\cdot\nabla\varphi\,dx
\qquad \forall\varphi\in H^1_0(D).
\]
Testing with \(\varphi=q\) gives
\[
\lambda\norm{\nabla q}_2^2
\le
\norm{g}_\infty\norm{\nabla u}_2\norm{\nabla q}_2.
\]
By Lemma~\ref{lem:forward-energy}, \(\norm{\nabla u}_2\le C_P/\lambda\), hence
\[
\norm{\nabla q}_2
\le
\frac{C_P}{\lambda^2}\norm{g}_\infty.
\]
By definition of \(v=A_{f'}g\),
\[
\int_D f'\nabla v\cdot\nabla\varphi\,dx
=
-\int_D g\nabla u'\cdot\nabla\varphi\,dx.
\]
Subtracting this from the equation for \(q\) yields
\[
\int_D f'\nabla R\cdot\nabla\varphi\,dx
=
-\int_D g\nabla q\cdot\nabla\varphi\,dx.
\]
Testing with \(\varphi=R\) gives
\[
\lambda\norm{\nabla R}_2^2
\le
\norm{g}_\infty\norm{\nabla q}_2\norm{\nabla R}_2.
\]
Therefore
\[
\norm{\nabla R}_2
\le
\frac1\lambda\norm{g}_\infty\norm{\nabla q}_2
\le
\frac{C_P}{\lambda^3}\norm{g}_\infty^2.
\]
Poincar\'e's inequality gives
\[
\norm{R}_2
\le
C_P\norm{\nabla R}_2
\le
\frac{C_P^2}{(1-r)^3}\norm{g}_\infty^2.
\]
Since \(\norm{g}_\infty\le\norm{g}_{C^\alpha}\), the theorem follows.
\end{proof}

\section{Frozen periodic linearized inverse problem}
\label{sec:frozen}

We now analyze a frozen periodic linearized model. The original PDE has variable coefficients and Dirichlet boundary conditions; the frozen model is a local periodic surrogate. We choose coordinates so that the frozen nonzero direction is \(e_d\). Let
\[
\T^d=\R^d/\Z^d.
\]
Fix \(\beta>0\) and define the Fourier multiplier
\[
a_\beta(\xi)=\frac{1+i\beta\xi_d}{|\xi|^2},
\qquad
\xi\in\Z^d\setminus\{0\}.
\]
Let \(A_\beta\) be the corresponding Fourier multiplier operator on mean-zero functions. The frozen Gaussian sequence model is
\begin{equation}
\label{eq:frozen-model}
Y_\xi=a_\beta(\xi)\theta_\xi+n^{-1/2}Z_\xi,
\qquad \xi\neq0,
\end{equation}
where \(\theta_\xi\) are the Fourier coefficients of a real mean-zero function \(\theta\), and \(Z_\xi\) denote the Fourier coordinates of Gaussian white noise. We suppress the usual conjugacy constraints for real signals; they affect constants only. The zero mode is set to zero: \(\theta_0=0\).

The parameter space is the Besov ball
\[
\Theta_\alpha(L)
=
\left\{
\theta\in B^\alpha_{\infty,\infty}(\T^d):
\norm{\theta}_{B^\alpha_{\infty,\infty}}\le L,
\ \theta_0=0
\right\},
\]
where \(\alpha>0\). We use the Littlewood--Paley characterization
\[
\norm{\theta}_{B^\alpha_{\infty,\infty}}
\asym
\sup_{j\ge0}2^{j\alpha}\norm{\Delta_j\theta}_{L^\infty};
\]
see \cite{Triebel}.

\subsection{Counting and variance}

% RECORDS: 199
\begin{lemma}[Weyl counting for the frozen symbol]
\label{lem:frozen-counting}
Fix \(\beta>0\). For \(0<\sigma\le\sigma_0(d,\beta)\), define
\[
N_\beta(\sigma)
=
\#\{\xi\in\Z^d\setminus\{0\}: |a_\beta(\xi)|>\sigma\}.
\]
Then
\[
N_\beta(\sigma)\asym_{d,\beta}\sigma^{-d}.
\]
\end{lemma}
\begin{proof}
For \(|\xi|\ge1\),
\[
|a_\beta(\xi)|
=
\frac{\sqrt{1+\beta^2\xi_d^2}}{|\xi|^2}
\le
C_\beta\frac{1+|\xi|}{|\xi|^2}
\le
\frac{C_\beta}{|\xi|}.
\]
Thus \(|a_\beta(\xi)|>\sigma\) implies \(|\xi|\le C_\beta/\sigma\), giving
\[
N_\beta(\sigma)\le C\sigma^{-d}.
\]

For the lower bound, restrict to the cone
\[
|\xi_d|\ge \frac12|\xi|.
\]
On this cone,
\[
|a_\beta(\xi)|
\ge
\frac{\beta|\xi_d|}{|\xi|^2}
\ge
\frac{\beta}{2|\xi|}.
\]
If
\[
1\le |\xi|\le c\beta/\sigma,
\]
then \(|a_\beta(\xi)|\ge\sigma\). The number of lattice points in this cone and ball is comparable to \(\sigma^{-d}\), with constants depending on \(d\) and \(\beta\). Hence
\[
N_\beta(\sigma)\ge c\sigma^{-d}.
\]
\end{proof}

The counting law suggests degree-one smoothing, but it is not the rate-determining quantity. The minimax rate is controlled by the variance sum.

% RECORDS: 199
\begin{lemma}[Frozen variance sum]
\label{lem:frozen-variance}
Fix \(\beta>0\). For \(R\ge1\),
\[
S_\beta(R)
=
\sum_{\substack{\xi\in\Z^d\\1\le|\xi|\le R}}
|a_\beta(\xi)|^{-2}
=
\sum_{\substack{\xi\in\Z^d\\1\le|\xi|\le R}}
\frac{|\xi|^4}{1+\beta^2\xi_d^2}
\asym_{d,\beta}
R^{d+3}.
\]
\end{lemma}
\begin{proof}
For the lower bound, take \(\xi_d=0\) and \(\xi'=(\xi_1,\dots,\xi_{d-1})\) in the annulus
\[
R/2\le |\xi'|\le R.
\]
There are \(\asym R^{d-1}\) such lattice points, and for each,
\[
\frac{|\xi|^4}{1+\beta^2\xi_d^2}
=
|\xi'|^4
\asym R^4.
\]
Thus
\[
S_\beta(R)\ge cR^{d+3}.
\]

For the upper bound, decompose into dyadic shells
\[
2^k\le |\xi|<2^{k+1},
\qquad
0\le k\le \lceil\log_2 R\rceil.
\]
On such a shell, \(|\xi|^4\le C2^{4k}\). Write \(\xi=(\eta,m)\) with \(\eta\in\Z^{d-1}\) and \(m=\xi_d\). For fixed \(\eta\) with \(|\eta|\le C2^k\),
\[
\sum_{m\in\Z}\frac{1}{1+\beta^2m^2}
\le
C_\beta.
\]
The number of such \(\eta\) is \(\asym 2^{k(d-1)}\). Hence the \(k\)-th shell contributes at most
\[
C_\beta 2^{4k}2^{k(d-1)}
=
C_\beta 2^{k(d+3)}.
\]
Summing over \(k\) up to \(\log_2 R\) gives
\[
S_\beta(R)\le C R^{d+3}.
\]
\end{proof}

% RECORDS: 199
\begin{lemma}[Besov projection tail]
\label{lem:besov-tail}
Let \(\theta\in\Theta_\alpha(L)\), and let \(\Pi_R\) be the Fourier projection onto \(\{|\xi|\le R\}\). Then
\[
\norm{\theta-\Pi_R\theta}_{L^2(\T^d)}^2
\le
C_{d,\alpha}L^2R^{-2\alpha}.
\]
\end{lemma}
\begin{proof}
Let \(\Delta_j\) be dyadic Littlewood--Paley projections. Choose \(J\) with \(2^J\le R<2^{J+1}\). The tail \(\theta-\Pi_R\theta\) is controlled by dyadic blocks with \(j\gtrsim J\). Since
\[
\norm{\Delta_j\theta}_{L^\infty}
\le
CL2^{-j\alpha},
\]
and \(\T^d\) has finite measure,
\[
\norm{\Delta_j\theta}_{L^2}
\le
CL2^{-j\alpha}.
\]
Therefore
\[
\norm{\theta-\Pi_R\theta}_{L^2}
\le
\sum_{j\ge J-1}\norm{\Delta_j\theta}_{L^2}
\le
CL\sum_{j\ge J-1}2^{-j\alpha}
\le
CL2^{-J\alpha}
\le
CLR^{-\alpha}.
\]
Squaring gives the claim.
\end{proof}

\subsection{Frozen minimax theorem}

% RECORDS: 199
\begin{lemma}[Wavelet packing in the equatorial slab]
\label{lem:wavelet-packing}
Let \(d'=d-1\). For every sufficiently large dyadic scale \(j\), there exist real orthonormal functions
\[
\psi_{j,1},\dots,\psi_{j,N_j}\in L^2(\T^{d'})
\]
with \(N_j\asym 2^{jd'}\), whose Fourier supports lie in a dyadic annulus
\[
c2^j\le |\eta|\le C2^j,
\]
and such that for any coefficients \(b_k\) with \(|b_k|\le b\),
\[
\left\|
\sum_{k=1}^{N_j}b_k\psi_{j,k}
\right\|_{L^\infty(\T^{d'})}
\le
C b\,2^{jd'/2}.
\]
Extending these functions constantly in the \(x_d\)-variable gives functions
\[
\Psi_{j,k}(x',x_d)=\psi_{j,k}(x')
\]
on \(\T^d\), orthonormal in \(L^2(\T^d)\), with Fourier support in
\[
\{(\eta,0): c2^j\le |\eta|\le C2^j\}.
\]
\end{lemma}
\begin{proof}
This is the standard periodized Meyer wavelet construction on the torus; see Meyer \cite[Chapter 3]{Meyer} and Daubechies \cite[Chapter 5]{Daubechies}. The periodized Meyer system is an orthonormal basis with Fourier support in dyadic annuli. At scale \(j\), the number of wavelets is comparable to \(2^{jd'}\), each has \(L^\infty\)-size \(O(2^{jd'/2})\), and the uniform overlap estimate
\[
\sum_k|\psi_{j,k}(x')|\le C2^{jd'/2}
\]
holds. The stated \(L^\infty\) bound follows immediately for bounded coefficients. Extension constantly in \(x_d\) preserves orthonormality and places all Fourier mass at \(\xi_d=0\).
\end{proof}

% RECORDS: 199
\begin{lemma}[Assouad hypercube lemma]
\label{lem:assouad}
Let \(\{\psi_k\}_{k=1}^N\) be an orthonormal system in a Hilbert space, and let
\[
\theta_\omega=\delta\sum_{k=1}^N\omega_k\psi_k,
\qquad
\omega\in\{\pm1\}^N.
\]
Let \(P_\omega\) be probability laws of an observation under \(\theta_\omega\). Suppose that whenever \(\omega,\omega'\) differ in exactly one coordinate,
\[
\KL(P_\omega\|P_{\omega'})\le\kappa.
\]
If \(\kappa\le1/8\), then for every estimator \(\widehat\theta\),
\[
\sup_{\omega}\E_\omega\norm{\widehat\theta-\theta_\omega}^2
\ge
cN\delta^2
\]
with an absolute constant \(c>0\).
\end{lemma}
\begin{proof}
Let \(\Omega\) be uniform on \(\{\pm1\}^N\), and let \(X\sim P_\Omega\). For any estimator \(\widehat\theta\), define bit estimators
\[
\widehat\Omega_k=\operatorname{sign}\langle\widehat\theta,\psi_k\rangle.
\]
Since \(\langle\theta_\Omega,\psi_k\rangle=\delta\Omega_k\), if \(\widehat\Omega_k\neq\Omega_k\), then
\[
|\langle\widehat\theta-\theta_\Omega,\psi_k\rangle|\ge\delta.
\]
By orthonormality,
\[
\norm{\widehat\theta-\theta_\Omega}^2
\ge
\delta^2\sum_{k=1}^N\mathbf 1_{\{\widehat\Omega_k\neq\Omega_k\}}.
\]
Averaging over \(\Omega\) gives
\[
\E\norm{\widehat\theta-\theta_\Omega}^2
\ge
\delta^2\sum_{k=1}^N\PP(\widehat\Omega_k\neq\Omega_k).
\]
For fixed \(k\), the average bit error is bounded below by the Bayes error for testing the two mixture laws \(P_{k,+}\) and \(P_{k,-}\) obtained by conditioning on \(\Omega_k=\pm1\). Thus
\[
\PP(\widehat\Omega_k\neq\Omega_k)
\ge
\frac12\left(1-\TV(P_{k,+},P_{k,-})\right).
\]
By convexity of KL divergence and the adjacent KL assumption,
\[
\KL(P_{k,+}\|P_{k,-})\le\kappa.
\]
Pinsker's inequality \cite{Tsybakov} gives
\[
\TV(P_{k,+},P_{k,-})
\le
\sqrt{\kappa/2}
\le
\frac12
\]
when \(\kappa\le1/8\). Hence each bit error is at least \(1/4\), and
\[
\E\norm{\widehat\theta-\theta_\Omega}^2
\ge
\frac14N\delta^2.
\]
The supremum risk is at least the average risk.
\end{proof}

The certified frozen statement is an up-to-logarithmic-factor exponent. The periodic wavelet construction actually gives the sharper constant-order bounds stated below.

% RECORDS: 199
\begin{theorem}[Frozen minimax rate]
\label{thm:frozen-minimax}
Fix \(\beta>0\), \(\alpha>0\), and \(L>0\). For the frozen periodic model \eqref{eq:frozen-model} over \(\Theta_\alpha(L)\), let
\[
\mathcal R_n^{\mathrm{fr}}
=
\inf_{\widehat\theta}\sup_{\theta\in\Theta_\alpha(L)}
\E\norm{\widehat\theta-\theta}_{L^2(\T^d)}^2.
\]
Then, with
\[
p=2\alpha+d+3,
\]
there exist constants \(c,C>0\), depending only on \(d,\alpha,\beta,L\), such that for all sufficiently large \(n\),
\[
c\,n^{-2\alpha/p}
\le
\mathcal R_n^{\mathrm{fr}}
\le
C\,n^{-2\alpha/p}.
\]
In particular, the frozen minimax exponent is \(2\alpha/(2\alpha+d+3)\); in the certified wording this is stated up to logarithmic factors, and the periodic wavelet construction removes the logarithm.
\end{theorem}
\begin{proof}
\textbf{Upper bound.} Fix \(R\ge1\) and define the sieve estimator
\[
\widehat\theta^R_\xi=
\begin{cases}
Y_\xi/a_\beta(\xi), & 1\le|\xi|\le R,\\
0, & \text{otherwise}.
\end{cases}
\]
Then
\[
\E\norm{\widehat\theta^R-\theta}_2^2
\le
C\left(
n^{-1}\sum_{1\le|\xi|\le R}|a_\beta(\xi)|^{-2}
+
\norm{\theta-\Pi_R\theta}_2^2
\right).
\]
By Lemma~\ref{lem:frozen-variance} and Lemma~\ref{lem:besov-tail},
\[
\sup_{\theta\in\Theta_\alpha(L)}
\E\norm{\widehat\theta^R-\theta}_2^2
\le
C\left(
\frac{R^{d+3}}{n}
+
L^2R^{-2\alpha}
\right).
\]
Choose
\[
R\asym (nL^2)^{1/(2\alpha+d+3)}.
\]
The two terms balance and yield
\[
\mathcal R_n^{\mathrm{fr}}
\le
C_{d,\alpha,\beta,L}
n^{-2\alpha/(2\alpha+d+3)}.
\]

\textbf{Lower bound.} Let \(d'=d-1\) and \(p=2\alpha+d+3=2\alpha+d'+4\). Choose a dyadic scale \(R=2^j\) comparable to \((nL^2)^{1/p}\), and let \(\{\Psi_k\}_{k=1}^N\) be the functions from Lemma~\ref{lem:wavelet-packing} at that scale, so \(N\asym R^{d'}\). Define
\[
\delta=c_0 L R^{-\alpha-d'/2},
\]
with \(c_0>0\) small. For \(\omega\in\{\pm1\}^N\), set
\[
\theta_\omega=\delta\sum_{k=1}^N\omega_k\Psi_k.
\]
The Fourier support of every \(\Psi_k\) lies in an annulus \(|\xi|\asym R\) with \(\xi_d=0\), so all \(\theta_\omega\) lie in a single Littlewood--Paley block. By Lemma~\ref{lem:wavelet-packing},
\[
\norm{\theta_\omega}_{L^\infty}
\le
C\delta R^{d'/2}
=
C c_0 L R^{-\alpha}.
\]
Since the block is at frequency \(R\),
\[
\norm{\theta_\omega}_{B^\alpha_{\infty,\infty}}
\le
C R^\alpha\norm{\theta_\omega}_{L^\infty}
\le
C c_0 L.
\]
Choosing \(c_0\) small enough gives \(\theta_\omega\in\Theta_\alpha(L)\).

For adjacent \(\omega,\omega'\) differing in one coordinate,
\[
\theta_\omega-\theta_{\omega'}=2\delta\Psi_k
\]
for some \(k\). Since \(\Psi_k\) has Fourier support with \(\xi_d=0\) and \(|\xi|\asym R\),
\[
|a_\beta(\xi)|^2
=
\frac1{|\xi|^4}
\asym R^{-4}
\]
on that support. Hence
\[
\norm{A_\beta(\theta_\omega-\theta_{\omega'})}_{L^2}^2
\le
C\delta^2R^{-4}.
\]
By Lemma~\ref{lem:white-kl}, the KL divergence between the corresponding Gaussian laws is
\[
\KL(P_\omega\|P_{\omega'})
=
\frac n2
\norm{A_\beta(\theta_\omega-\theta_{\omega'})}_2^2
\le
C n\delta^2R^{-4}.
\]
Using \(R^p\asym nL^2\) and \(\delta^2=c_0^2L^2R^{-2\alpha-d'}\), we obtain
\[
n\delta^2R^{-4}
\asym
c_0^2 nL^2R^{-p}
\asym
c_0^2.
\]
Choosing \(c_0\) small makes this at most \(1/8\). Lemma~\ref{lem:assouad} gives
\[
\mathcal R_n^{\mathrm{fr}}
\ge
cN\delta^2
\asym
R^{d'}L^2R^{-2\alpha-d'}
=
L^2R^{-2\alpha}.
\]
Since \(R\asym (nL^2)^{1/p}\),
\[
L^2R^{-2\alpha}
\asym
n^{-2\alpha/p}.
\]
This proves the lower bound.
\end{proof}

\section{Variable-coefficient phase-space variance}
\label{sec:variable-variance}

For the true linearization at \(f=1\), freezing at \(x\) gives the symbol
\[
a(x,\xi)=\frac{1+i(x/d)\cdot\xi}{|\xi|^2}.
\]
The corresponding phase-space variance is
\[
S_D(R)
=
\int_D\int_{|\xi|\le R}
|a(x,\xi)|^{-2}\,d\xi\,dx
=
\int_D\int_{|\xi|\le R}
\frac{|\xi|^4}{1+|(x/d)\cdot\xi|^2}\,d\xi\,dx.
\]

% RECORDS: 199
\begin{lemma}[Angular integral]
\label{lem:angular-integral}
Let \(e\in S^{d-1}\) and \(t\ge0\). Define
\[
J(t)=\int_{S^{d-1}}\frac{d\sigma(\omega)}{1+t^2(\omega\cdot e)^2}.
\]
Then
\[
J(t)\asym_d \min\{1,t^{-1}\},
\]
with the convention \(t^{-1}=\infty\) at \(t=0\).
\end{lemma}
\begin{proof}
Write \(s=\omega\cdot e\). Then
\[
J(t)=c_d\int_{-1}^1
\frac{(1-s^2)^{(d-3)/2}}{1+t^2s^2}\,ds
\]
for \(d\ge3\), with the usual integrable endpoint weight for \(d=2\).

If \(0\le t\le1\), then \(1\le1+t^2s^2\le2\), so \(J(t)\asym1\).

Let \(t\ge1\). For the lower bound, restrict to \(|s|\le1/(2t)\). Near \(s=0\), the weight is bounded below and the denominator is at most \(2\). Hence
\[
J(t)\ge c\int_{-1/(2t)}^{1/(2t)}ds=\frac{c}{t}.
\]
For the upper bound, split into \(|s|\le1/t\) and \(|s|>1/t\). On \(|s|\le1/t\), the weight is integrable near \(0\), giving a contribution \(O(1/t)\). On \(|s|>1/t\), the denominator is at least \(t^2s^2\), so the contribution is bounded by
\[
t^{-2}\int_{|s|>1/t}
\frac{(1-s^2)^{(d-3)/2}}{s^2}\,ds.
\]
For \(d\ge3\), the weight is bounded, and
\[
\int_{1/t}^1s^{-2}\,ds\le Ct.
\]
For \(d=2\), the weight is \((1-s^2)^{-1/2}\); near \(s=0\) the same bound gives \(O(t)\), while near the endpoints the integral is \(O(1)\). Hence in all cases the second contribution is \(O(1/t)\). Therefore \(J(t)\le C/t\).
\end{proof}

% RECORDS: 199,44
\begin{theorem}[Variable-coefficient variance order]
\label{thm:phase-space-variance}
For \(d\ge2\) and \(R\ge1\),
\[
S_D(R)\asym_d R^{d+3}.
\]
\end{theorem}
\begin{proof}
For fixed \(x\), put \(\rho=|x|/d\). By rotational symmetry of the \(\xi\)-integral,
\[
\int_{|\xi|\le R}
\frac{|\xi|^4}{1+|(x/d)\cdot\xi|^2}\,d\xi
=
\int_0^R r^{d+3}J(\rho r)\,dr.
\]
By Lemma~\ref{lem:angular-integral},
\[
J(\rho r)\asym
\begin{cases}
1, & \rho r\le1,\\
(\rho r)^{-1}, & \rho r\ge1.
\end{cases}
\]
If \(R\rho\le1\), the integral is \(\asym R^{d+4}\). If \(R\rho>1\), the dominant contribution is
\[
\asym
\rho^{-1}\int_{1/\rho}^R r^{d+2}\,dr
\asym
\rho^{-1}R^{d+3}.
\]

Split \(D\) into
\[
D_0=\{x: |x|\le d/R\},
\qquad
D_1=D\setminus D_0.
\]
On \(D_0\), the inner integral is \(O(R^{d+4})\), while \(|D_0|\asym R^{-d}\); the contribution is \(O(R^4)\), which is lower order because \(d\ge2\).

On \(D_1\), the contribution is comparable to
\[
R^{d+3}\int_{D_1}\frac{d}{|x|}\,dx.
\]
Since \(d\ge2\),
\[
\int_{d/R<|x|<1}\frac{dx}{|x|}
=
|S^{d-1}|\int_{d/R}^1 r^{d-2}\,dr
\asym1.
\]
Thus the outer contribution is \(\asym R^{d+3}\). The lower bound also follows directly by restricting to \(1/2\le|x|\le1\), where \(\rho\) is bounded below. Hence
\[
S_D(R)\asym R^{d+3}.
\]
\end{proof}

This variance order matches the frozen-model order. It supports the candidate exponent \(2\alpha/(2\alpha+d+3)\) for the original linearized operator, but by itself it is neither an estimator upper bound nor a lower-bound construction for the full variable-coefficient nonlinear problem.

\section{Critical-point two-point lower bound}
\label{sec:critical}

At the origin, \(\nabla u_1(0)=0\), and the linearized operator degenerates further. We prove a rigorous two-point Le Cam lower bound based on origin-localized radial bumps.

Fix a nonzero radial \(\psi\in C_c^\infty(B(0,1))\) with \(\psi\ge0\) and \(\norm{\psi}_{L^2}=1\). For \(a>0\) and \(0<\varepsilon<1/2\), define
\[
h_{\varepsilon,a}(x)=a\psi(x/\varepsilon).
\]

% RECORDS: 756
\begin{lemma}[Degree-two gain at the critical point]
\label{lem:critical-gain}
Let \(A=A_1\). For radial \(\psi\) as above,
\[
\norm{A h_{\varepsilon,a}}_{L^2(D)}
=
a\varepsilon^{2+d/2}\bigl(C_\psi+o(1)\bigr)
\]
as \(\varepsilon\downarrow0\), where \(C_\psi\in(0,\infty)\) depends only on \(\psi\). Moreover,
\[
\norm{\nabla A h_{\varepsilon,a}}_{L^2(D)}
\le
C_\psi' a\varepsilon^{1+d/2}.
\]
\end{lemma}
\begin{proof}
Set \(y=x/\varepsilon\). Then
\[
h_{\varepsilon,a}(x)+\frac{x}{d}\cdot\nabla h_{\varepsilon,a}(x)
=
a\left(
\psi(y)+\frac{y\cdot\nabla\psi(y)}{d}
\right)
=:aS(y).
\]
But
\[
S(y)=\frac1d\divg_y(y\psi(y)).
\]
Let \(\widehat w\) solve on \(\R^d\)
\[
-\Delta_y\widehat w=S,
\]
with decay at infinity. Since \(\psi\) is radial, \(S\) is radial, compactly supported in \(B(0,1)\), and has mean zero:
\[
\int_{\R^d}S(y)\,dy
=
\frac1d\int_{\R^d}\divg(y\psi)\,dy
=
0.
\]
For \(|y|>1\), the radial solution satisfies
\[
\widehat w'(r)
=
-r^{1-d}\int_0^r s^{d-1}S(s)\,ds.
\]
For \(r\ge1\), the integral is the total mass of \(S\), hence zero. Thus \(\widehat w'=0\) for \(r\ge1\), and decay forces \(\widehat w=0\) there. Therefore \(\widehat w\) is compactly supported in \(B(0,1)\).

Define
\[
w(x)=a\varepsilon^2\widehat w(x/\varepsilon).
\]
Then
\[
\Delta_x w
=
a\Delta_y\widehat w(y)
=
-aS(y),
\]
so \(w\) solves the Dirichlet problem for \(Ah_{\varepsilon,a}\), because its support lies in \(B(0,\varepsilon)\subset D\). Hence
\[
\norm{Ah_{\varepsilon,a}}_2^2
=
a^2\varepsilon^4
\int_D|\widehat w(x/\varepsilon)|^2\,dx
=
a^2\varepsilon^{4+d}\norm{\widehat w}_2^2.
\]
Set \(C_\psi=\norm{\widehat w}_2\). If \(\widehat w=0\), then \(S=0\), hence \(\divg(y\psi)=0\). For smooth compactly supported radial \(\psi\), this forces \(\psi=0\), contradicting the choice of \(\psi\). Thus \(C_\psi>0\).

Finally,
\[
\nabla w(x)=a\varepsilon\nabla\widehat w(x/\varepsilon),
\]
so
\[
\norm{\nabla Ah_{\varepsilon,a}}_2
=
a\varepsilon^{1+d/2}\norm{\nabla\widehat w}_2,
\]
which gives the gradient bound.
\end{proof}

% RECORDS: 756
\begin{lemma}[Nonlinear remainder near the critical point]
\label{lem:critical-remainder}
Let \(u_h=\Phi(1+h)\) with \(h=h_{\varepsilon,a}\) as above, and let \(w_1=Ah\). Assume \(a\) is small enough that \(1+h\ge1/2\). Then the remainder \(w_2=u_h-u_1-w_1\) satisfies
\[
\norm{w_2}_{L^2(D)}
\le
C a^2\varepsilon^{1+d/2}.
\]
Consequently, if \(a=c_0r\varepsilon^\alpha\) with \(\alpha>1\), then
\[
\frac{\norm{w_2}_{L^2}}{\norm{Ah}_{L^2}}
\le
C c_0r\,\varepsilon^{\alpha-1}\to0.
\]
\end{lemma}
\begin{proof}
Let \(w=u_h-u_1\). Since \(u_h\) solves \(\divg((1+h)\nabla u_h)=1\) and \(u_1\) solves \(\Delta u_1=1\),
\[
\divg((1+h)\nabla w)+\divg(h\nabla u_1)=0.
\]
Weakly,
\[
\int_D(1+h)\nabla w\cdot\nabla\varphi
=
-\int_D h\nabla u_1\cdot\nabla\varphi.
\]
The linear part \(w_1=Ah\) satisfies
\[
\int_D\nabla w_1\cdot\nabla\varphi
=
-\int_D h\nabla u_1\cdot\nabla\varphi.
\]
Subtracting, \(w_2=w-w_1\) satisfies
\[
\int_D(1+h)\nabla w_2\cdot\nabla\varphi
=
-\int_D h\nabla w_1\cdot\nabla\varphi.
\]
Testing with \(\varphi=w_2\) and using \(1+h\ge1/2\) gives
\[
\norm{\nabla w_2}_2
\le
C\norm{h\nabla w_1}_2
\le
C\norm{h}_\infty\norm{\nabla w_1}_2.
\]
By Lemma~\ref{lem:critical-gain},
\[
\norm{\nabla w_1}_2
\le
C a\varepsilon^{1+d/2},
\]
and \(\norm{h}_\infty\le a\). Therefore
\[
\norm{\nabla w_2}_2
\le
C a^2\varepsilon^{1+d/2}.
\]
Poincar\'e's inequality gives the \(L^2\) bound.

By Lemma~\ref{lem:critical-gain},
\[
\norm{Ah}_2\asym a\varepsilon^{2+d/2}.
\]
Hence
\[
\frac{\norm{w_2}_2}{\norm{Ah}_2}
\le
C a\varepsilon^{-1}
=
C c_0r\varepsilon^{\alpha-1},
\]
which tends to zero for \(\alpha>1\).
\end{proof}

% RECORDS: 756
\begin{lemma}[Two-point Le Cam bound]
\label{lem:lecam}
Let \(P_0,P_1\) be probability measures with \(\KL(P_0\|P_1)\le\kappa\le1/8\). Let \(f_0,f_1\) be two parameters with separation \(\Delta=\norm{f_1-f_0}^2\). Then
\[
\inf_{\widehat f}\max_{i=0,1}\E_i\norm{\widehat f-f_i}^2
\ge
\frac{\Delta}{16}.
\]
\end{lemma}
\begin{proof}
For any estimator \(\widehat f\), define a test by choosing the closer parameter:
\[
\varphi=
\begin{cases}
0, & \norm{\widehat f-f_0}\le\norm{\widehat f-f_1},\\
1, & \text{otherwise}.
\end{cases}
\]
If \(f_0\) is true and \(\varphi=1\), then by the triangle inequality,
\[
\sqrt{\Delta}
\le
\norm{\widehat f-f_0}+\norm{\widehat f-f_1}
\le
2\norm{\widehat f-f_0},
\]
so \(\norm{\widehat f-f_0}^2\ge\Delta/4\). Similarly under \(f_1\), if \(\varphi=0\), then \(\norm{\widehat f-f_1}^2\ge\Delta/4\). Therefore
\[
\max_i\E_i\norm{\widehat f-f_i}^2
\ge
\frac{\Delta}{8}\bigl(1-\TV(P_0,P_1)\bigr).
\]
Pinsker's inequality \cite{Tsybakov} gives
\[
\TV(P_0,P_1)\le\sqrt{\kappa/2}\le\frac12,
\]
and the claim follows.
\end{proof}

% RECORDS: 756
\begin{theorem}[Critical-point lower bound for the original problem]
\label{thm:critical-lower}
Let \(d\ge2\), \(\alpha>1\), and fix \(r\in(0,1)\). There exists \(c>0\) such that for all large \(n\),
\[
\inf_{\widehat f_n}\sup_{f\in\mathcal F(\alpha,r)}
\E_f\norm{\widehat f_n-f}_{L^2(D)}^2
\ge
c\,n^{-\frac{2\alpha+d}{2\alpha+d+4}}.
\]
\end{theorem}
\begin{proof}
Let
\[
q=2\alpha+d+4.
\]
Choose \(c_0>0\) small enough that
\[
c_0\norm{\psi}_{C^\alpha}\le \frac14,
\qquad
c_0\norm{\psi}_\infty\le \frac14.
\]
Define
\[
\varepsilon_n=(c_1nr^2)^{-1/q}
\]
with \(c_1>0\) large, and set
\[
a_n=c_0r\varepsilon_n^\alpha,
\qquad
f_1=1+a_n\psi(\cdot/\varepsilon_n),
\qquad
f_0=1.
\]
By the standard scaling of H\"older--Zygmund norms,
\[
\norm{f_1-1}_{C^\alpha}
\le
C c_0r
\le r
\]
for \(c_0\) small, so \(f_1\in\mathcal F(\alpha,r)\). Also \(f_0\in\mathcal F(\alpha,r)\), and \(f_1\ge1\) because \(\psi\ge0\).

By Lemma~\ref{lem:critical-gain} and Lemma~\ref{lem:critical-remainder}, for large \(n\),
\[
\norm{u_{f_1}-u_{f_0}}_{L^2}
\asym
a_n\varepsilon_n^{2+d/2}
=
c_0r\varepsilon_n^{\alpha+2+d/2}.
\]
Thus
\[
\norm{u_{f_1}-u_{f_0}}_2^2
\asym
r^2\varepsilon_n^{2\alpha+d+4}.
\]
By Lemma~\ref{lem:white-kl}, the KL divergence between the white-noise laws is
\[
\KL(P_0\|P_1)
=
\frac n2\norm{u_{f_1}-u_{f_0}}_2^2
\asym
nr^2\varepsilon_n^{2\alpha+d+4}.
\]
By the choice of \(\varepsilon_n\), this is bounded by a universal constant; choosing \(c_1\) large makes it at most \(1/8\).

The parameter separation is
\[
\norm{f_1-f_0}_2^2
=
a_n^2\varepsilon_n^d\norm{\psi}_2^2
=
c_0^2r^2\varepsilon_n^{2\alpha+d}.
\]
Substituting \(\varepsilon_n=(c_1nr^2)^{-1/q}\) gives
\[
\norm{f_1-f_0}_2^2
\asym
r^2(nr^2)^{-\frac{2\alpha+d}{2\alpha+d+4}}.
\]
For fixed \(r\), this is comparable to
\[
n^{-\frac{2\alpha+d}{2\alpha+d+4}}.
\]
Lemma~\ref{lem:lecam} yields the lower bound.
\end{proof}

% RECORDS: 756,199
\begin{proposition}[Comparison of exponents]
\label{prop:exponent-comparison}
For \(d\ge2\) and \(\alpha>0\), define
\[
e_{\mathrm{crit}}=\frac{2\alpha+d}{2\alpha+d+4},
\qquad
e_{\mathrm{tan}}=\frac{2\alpha}{2\alpha+d+3},
\qquad
e_{\mathrm{cand}}=\frac{2\alpha}{2\alpha+d+2}.
\]
Then
\[
e_{\mathrm{tan}}<e_{\mathrm{cand}}<e_{\mathrm{crit}}.
\]
Thus the critical-point lower bound is strictly weaker than the frozen/tangential exponent, and the frozen exponent is strictly worse than the naive degree-one candidate.
\end{proposition}
\begin{proof}
First,
\[
e_{\mathrm{cand}}>e_{\mathrm{tan}}
\]
is immediate because \(2\alpha+d+2<2\alpha+d+3\).

Next, with \(A=2\alpha+d\),
\[
e_{\mathrm{crit}}-e_{\mathrm{tan}}
=
\frac{A}{A+4}-\frac{2\alpha}{A+3}
=
\frac{A(A+3)-2\alpha(A+4)}{(A+4)(A+3)}.
\]
The numerator is
\[
(2\alpha+d)(2\alpha+d+3)-2\alpha(2\alpha+d+4)
=
2\alpha(d-1)+d(d+3)>0
\]
for \(d\ge2\).

Finally,
\[
e_{\mathrm{crit}}-e_{\mathrm{cand}}
=
\frac{A}{A+4}-\frac{2\alpha}{A+2}
=
\frac{A(A+2)-2\alpha(A+4)}{(A+4)(A+2)}.
\]
The numerator is
\[
(2\alpha+d)(2\alpha+d+2)-2\alpha(2\alpha+d+4)
=
2\alpha(d-2)+d(d+2)>0
\]
for \(d\ge2\). This proves the ordering.
\end{proof}

\section{Open points and omitted material}
\label{sec:open}

The original minimax rate over \(\mathcal F(\alpha,r)\) remains open. The certified results above identify the candidate exponent
\[
\frac{2\alpha}{2\alpha+d+3}
\]
from the frozen linearized model and give supporting phase-space variance for the variable-coefficient linearization, but they do not prove the full variable-coefficient nonlinear minimax theorem. The missing ingredients include:
\begin{enumerate}[label=(\arabic*)]
\item a variable-coefficient upper bound achieving the candidate exponent, using a sieve or spectral estimator together with the moving-anchor stability of Theorem~\ref{thm:moving-anchor};
\item a matching lower bound for the original problem using tangential needlets with vanishing moments and a \(C^\alpha\)-admissible taper away from the origin;
\item rigorous transfer from the frozen periodic multiplier model to the Dirichlet ball operator, including boundary effects and spatial localization;
\item extension beyond the critical-point two-point lower bound of Theorem~\ref{thm:critical-lower}.
\end{enumerate}
The frozen exponent is certified up to logarithmic factors; the periodic wavelet construction gives a sharper constant-order statement only for the frozen surrogate. A Bayesian posterior-contraction direction, numerical smoothing checks, and schema/connectivity probes are omitted as unrelated to the certified core.

\begin{thebibliography}{99}

\bibitem{Evans}
L.~C.~Evans,
\emph{Partial Differential Equations}, 2nd ed.,
American Mathematical Society, Providence, RI, 2010.

\bibitem{Triebel}
H.~Triebel,
\emph{Theory of Function Spaces},
Birkh\"auser, Basel, 1983.

\bibitem{Meyer}
Y.~Meyer,
\emph{Wavelets and Operators},
Cambridge University Press, Cambridge, 1992.

\bibitem{Daubechies}
I.~Daubechies,
\emph{Ten Lectures on Wavelets},
SIAM, Philadelphia, 1992.

\bibitem{Tsybakov}
A.~B.~Tsybakov,
\emph{Introduction to Nonparametric Estimation},
Springer, New York, 2009.

\end{thebibliography}

\end{document}